
\medskip

\centerline{\large\bf Quot Scheme in Quasi-Projective case}


\medskip

\begin{claim}[Support-disjointness open locus]\label{nhq-support-disjoint-open}\source{nitsure-hilbert-quot:page-0028}
Let $\pi: Z\to S$ be a proper morphism of noetherian schemes, let
$Y\subset Z$ be a closed subscheme, and let $\F$ be a coherent sheaf on
$Z$. Then there exists an open subscheme $S'\subset S$ such that a
morphism $T\to S$ factors through $S'$ if and only if the support of
$\F_T$ on $Z_T$ is disjoint from $Y_T$.
\end{claim}

\begin{proof}
\notready The source poses this as an exercise and gives no proof.
\end{proof}

\begin{claim}[Open Quot subfunctor]\label{nhq-open-quot-subfunctor}\source{nitsure-hilbert-quot:page-0028}\uses{nhq-support-disjoint-open}
If $\pi: Z\to S$ is proper with $S$ noetherian, $X\subset Z$ is open,
and $E$ is coherent on $Z$, then $\Quot_{E|_X/X/S}$ is an open subfunctor
of $\Quot_{E/Z/S}$.
\end{claim}

\begin{proof}
\notready The source derives this consequence as an exercise from the preceding claim.
\end{proof}

With the preceding construction, a quot scheme extends to the strongly
quasi-projective case, to give the following.

\begin{theorem}\label{nhq-strongly-quasi-projective-quot}\dcref{6.1}\source{nitsure-hilbert-quot:page-0028}\uses{nhq-altman-kleiman-quot,nhq-open-quot-subfunctor,nhq-relative-quot-family}
{\rm (Altman and Kleiman)}
Let $S$ be a noetherian scheme, $X$ a locally closed subscheme of
$\PP(V)$ for some vector bundle $V$ on $S$, $L = {\mathcal O}_{\PP(V)}(1)|_X$,
$E$ a coherent quotient sheaf of $\pi^*(W)(\nu)|_X$ where
$W$ is a vector bundle on $S$ and $\nu$ is an integer, and
$\Phi \in \Q[\lambda]$.
Then the functor $\Quot_{E/X/S}^{\Phi,L}$ is representable
by a scheme $\quot_{E/X/S}^{\Phi,L}$ which can be embedded over $S$ as a
locally closed subscheme of $\PP(F)$ for some vector bundle $F$ on $S$.
Moreover, the vector bundle $F$ can be chosen to be an exterior power of
the tensor product of $W$ with a symmetric power of $V$.
\end{theorem}

\begin{proof} Let $\ov{X} \subset \PP(V)$ be the schematic closure of
$X\subset \PP(V)$, and let $\ov{E}$ be the coherent sheaf on $\ov{X}$
defined as the image of the composite homomorphism
$$\pi^*(W)(\nu)|_{\ov{X}} \to j_*(\pi^*(W)(\nu)|_X) \to j_*E$$
Then we get a quotient $\pi^*(W)(\nu)|_{\ov{X}} \to \ov{E}$
which restricts on $X\subset \ov{X}$ to the given quotient
$\pi^*(W)(\nu)|_X\to E$. The condition that the quotient remain supported
on $X$ defines an open subfunctor, so
$\Quot_{E/X/S}$ is an open subfunctor of
$\Quot_{\ov{E}/\ov{X}/S}$. Now the result follows from the
Theorem \ref{nhq-altman-kleiman-quot}.\end{proof}



\medskip

In order to extend Grothendieck's construction of a quot
scheme to the quasi-projective case, one first needs the following lemma
which is of independent interest.


\begin{lemma}\label{nhq-coherent-prolongation}\dcref{6.2}\source{nitsure-hilbert-quot:page-0028}
Any coherent sheaf on an open subscheme of a noetherian scheme $S$ can be
prolonged to a coherent sheaf on all of $S$.
\end{lemma}

\begin{proof} First consider the case where $S = \spec A$ is affine, and
let $j: U\hra S$ denote the inclusion.
The quasi-coherent sheaf
$j_*(\F)$ corresponds to the $A$-module $M = H^0(S,j_*(\F))$,
in the sense that $j_*(\F) = M^{\sim}$.
Given any $u\in U$, there exist finitely many
elements $e_1,\ldots,e_n \in M$
which generate the fiber $\F_u$ regarded as a vector space
over the residue field $\kappa(u)$. By Nakayama these elements
will generate the stalks of $\F$ in an open neighbourhood of $u$ in $U$.
Therefore by the noetherian hypothesis, there exist finitely many
elements $e_1,\ldots,e_r \in M$ which generate the stalk of $\F$
at each point of $U$. If $N\subset M$ is the submodule generated by these
elements, then $\G = N^{\sim}$ is a coherent prolongation of
$\F$ to $S = \spec A$, proving the result in the affine case.

In the general case, by the
noetherian condition there exists a maximal coherent prolongation
$(U',\F')$ of $\F$. Then unless $U' = S$, we can obtain
a further prolongation of $\F'$ by using the affine case.
For, if $u\in S- U'$,
we can take an affine open subscheme $V$ containing $u$, and a
coherent prolongation $\G'$ of $\F'|_{U'\bigcap V}$ to all of $V$,
and then glue together $\G'$ and $\F'$ along $U'\bigcap V$ to
further prolong $\F'$ to $U'\bigcup V$,
contradicting the maximality of $(U',\F')$.
\end{proof}


\begin{theorem}\label{nhq-quasi-projective-quot}\dcref{6.3}\source{nitsure-hilbert-quot:page-0029}\uses{nhq-coherent-prolongation,nhq-grothendieck-quot}{\rm (Grothendieck)}
Let $S$ be a noetherian scheme, $X$ a quasi-projective
scheme over $S$, $L$ a line bundle on $X$ which is
relatively very ample over $S$, $E$ a quotient sheaf on $X$,
and $\Phi \in \Q[\lambda]$.
Then the functor $\Quot_{E/X/S}^{\Phi,L}$ is representable
by a scheme $\quot_{E/X/S}^{\Phi,L}$ which is
quasi-projective over $S$.
\end{theorem}

\begin{proof} By definition of quasi-projectivity of $X\to S$, note that
$X$ can be embedded over $S$ as a locally closed subscheme of
$\PP(V)$ for some coherent sheaf $V$ on $S$,
such that $L$ is isomorphic to ${\mathcal O}_{\PP(V)}(1)|_X$.
Let $\ov{X} \subset \PP(V)$ be the schematic closure of
$X$ in $\PP(V)$. This is a projective scheme over $S$, and
$X$ is embedded as an open subscheme in it.
By Lemma \ref{nhq-coherent-prolongation} the
coherent sheaf $E$ has a coherent prolongation $\ov{E}$ to $\ov{X}$.
For any such prolongation $\ov{E}$, the functor
$\Quot_{E/X/S}$ is an open subfunctor of $\Quot_{\ov{E}/\ov{X}/S}$.
Therefore the desired result now follows from Theorem \ref{nhq-grothendieck-quot}.
\end{proof}







\bigskip


\centerline{\large\bf Scheme of Morphisms}


\medskip


We recall the following basic facts about flatness.

\begin{lemma}\label{nhq-local-criterion}\dcref{6.4}\source{nitsure-hilbert-quot:page-0029}
\textbf{(1) } Any finite-type flat morphism between noetherian schemes is open.


\textbf{(2) } Let $\pi : Y \to X$ be a finite-type morphism of
noetherian schemes. Then all $y\in Y$ such that $\pi$ is
flat at $y$ (that is, ${\mathcal O}_{Y,y}$ is a flat ${\mathcal O}_{X,\pi(y)}$-module)
form an open subset of $Y$.


\textbf{(3) } Let $S$ be a noetherian scheme, and let $f:X \to S$
and $g: Y\to S$ be finite type flat morphisms.
Let $\pi : Y \to X$ be any morphism such that $g = f\circ \pi$.
Let $y\in Y$, let $x=\pi(y)$, and let
$s = g(y) = f(x)$. If the restricted
morphism $\pi_s : Y_s \to X_s$ between the
fibers over $s$ is flat at $y\in Y_s$, then $\pi$ is flat at $y\in Y$.
\end{lemma}

\begin{proof} See for example Altman and Kleiman [A-K 1] Chapter V.
The statement \textbf{(3) } is a consequence of what is known as the
\textbf{local criterion for flatness}. \notready\end{proof}


\begin{theorem}\label{nhq-isomorphism-open}\dcref{6.5}\source{nitsure-hilbert-quot:page-0030}\uses{nhq-local-criterion}
Let $S$ be a noetherian scheme, and let $f:X \to S$
and $g: Y\to S$ be proper flat morphisms.
Let $\pi : Y \to X$ be any projective morphism with
$g = f\circ \pi$. Then $S$ has open subschemes
$S_2\subset S_1 \subset S$ with the following universal properties:

\textbf{(a) } For any locally noetherian $S$-scheme $T$, the base change
$\pi_T : Y_T \to X_T$ is a flat morphism if and only if the
structure morphism $T\to S$ factors via $S_1$. (This does not need
$\pi$ to be projective.)


\textbf{(b) } For any locally noetherian $S$-scheme $T$, the base change
$\pi_T : Y_T \to X_T$ is an isomorphism if and only if the
structure morphism $T\to S$ factors via $S_2$.
\end{theorem}

\begin{proof} \textbf{(a) } By Lemma \ref{nhq-local-criterion}.(2), all
$y\in Y$ such that $\pi$ is flat at $y$ form an
open subset $Y'\subset Y$. Then $S_1 = S - g(Y - Y')$
is an open subset of $S$ as $g$ is proper.
We give $S_1$ the open subscheme structure induced from $S$.
It follows from the local criterion of flatness
(Lemma  \ref{nhq-local-criterion}.(3))
that $S_1$ exactly consists of all $s\in S$ such that
the restricted morphism $\pi_s : Y_s \to X_s$ between the
fibers over $s$ is flat. Therefore again by the
local criterion of flatness, $S_1$ has the desired universal property.

\textbf{(b) } Let $\pi_1 : Y_1 \to X_1$ be the pull-back of
$\pi$ under the inclusion $S_1 \hra S$.
Let $L$ be a relatively very ample
line bundle for the projective morphism $\pi_1 : Y_1\to X_1$.
Then by noetherianness there exists an integer $m\ge 1$ such that
${\pi_1}_*L^m$ is generated by its global sections and $R^i{\pi_1}_*L^m =0$
for all $i \ge 1$. By flatness of $\pi_1$, it follows that
${\pi_1}_*L^m$ is a locally free sheaf. Let $U\subset X_1$ be the
open subschemes such that
${\pi_1}_*L^m$ is of rank $1$ on $U$.
Finally, let $S_2 = S_1 - f(X_1 - U)$, which is open as
$f$ is proper.
We give $S_2$ the induced open subscheme structure,
and leave it to the reader to verify that
it indeed has the required universal property \textbf{(b) }. \notready\end{proof}




\medskip

\begin{definition}[Parametrized $S$-morphisms]\label{nhq-parametrized-s-morphism}\source{nitsure-hilbert-quot:page-0031}
If $X$ and $Y$ are schemes over a base $S$, then for any $S$-scheme $T$,
an \textbf{$S$-morphism from $X$ to $Y$ parametrised by $T$} means a
$T$-morphism from $X\times_ST$ to $Y\times_ST$. The set of all such
morphisms is denoted by $\Mor_S(X,Y)(T)$, and
$T\mapsto\Mor_S(X,Y)(T)$ defines a contravariant functor
$\Mor_S(X,Y)$ from $S$-schemes to $Sets$.
\end{definition}


\medskip



\begin{theorem}\label{nhq-morphism-scheme}\dcref{6.6}\source{nitsure-hilbert-quot:page-0031}\uses{nhq-parametrized-s-morphism,nhq-isomorphism-open,nhq-hilbert-functor}
Let $S$ be a noetherian scheme, let $X$ be a projective scheme over $S$,
and let $Y$ be quasi-projective scheme
over $S$. Assume moreover that $X$ is flat over $S$. Then the
functor $\Mor_S(X,Y)$ is representable by an open
subscheme $\mor_S(X,Y)$ of $\hilb_{X\times_SY/S}$.
\end{theorem}

\begin{proof} We can associate
to each morphism $f : X_T \to Y_T$ (where $T$ is a scheme over $S$)
its graph $\Gamma_T(f) \subset (X\times_SY)_T$, which is
closed in $(X\times_SY)_T$ by separatedness of $Y\to S$.
We regard $\Gamma_T(f)$ as a closed subscheme of $(X\times_SY)_T$
which is isomorphic to $X$ under the graph morphism
$(\id_X,f): X \to  \Gamma_T(f)$ and the projection $\Gamma_T(f) \to X$,
which are inverses to each other.
As $X$ is proper and flat over $S$, so is $\Gamma_T(f)$, therefore
this defines a set-map $\Gamma_T : \Mor_S(X,Y)(T) \to
\Hilb_{X\times_SY/S}(T)$ which is functorial in $T$, so
we obtain a morphism of functors
$$\Gamma : \Mor_S(X,Y) \to \Hilb_{X\times_SY/S}$$


Given any element of $\Hilb_{X\times_SY/S}(T)$, represented by a
family $Z \subset (X\times_SY)_T$, it follows by applying
Theorem \ref{nhq-isomorphism-open}.(b) to
the projection $Z\to X$ that $T$ has an open subscheme
$T'$ with the following universal property: for any
base-change $U \to T$, the pull-back $Z_U\subset (X\times_SY)_U$
maps isomorphically on to $X_U$ under the projection
$p : (X\times_SY)_U\to X_U$ if and only if $U\to T$ factors via
$T'$. Note therefore that over $T'$,
the scheme $Z_{T'}$ will be the graph of a uniquely determined
morphism $X_{T'} \to Y_{T'}$.

This shows that the morphism of functors
$\Gamma : \Mor_S(X,Y) \to \Hilb_{X\times_SY/S}$
is a representable morphism which is an open embedding.
Therefore a representing scheme
$\mor_S(X,Y)$ for $\Mor_S(X,Y)$ exists as an open
subscheme of $\hilb_{X\times_SY/S}$. \end{proof}




\centerline{\large\bf Quotient by a Flat Projective Equivalence Relation}



\medskip


Let $X$ be a scheme over a base $S$.
\begin{definition}[Schematic equivalence relation]\label{nhq-schematic-equivalence-relation}\source{nitsure-hilbert-quot:page-0031}
A \textbf{schematic equivalence relation} on $X$ over $S$ will
mean an $S$-scheme $R$ together with a morphism
$f: R \to X\times_SX$ over $S$ such that for any $S$-scheme
$T$ the set map $f(T) : R(T) \to X(T)\times X(T)$
is injective and its image is the graph of an equivalence relation on
$X(T)$. (Here, we denote by $Z(T)$ the set $Mor_S(T,Z)= h_Z(T)$ of
all $T$-valued points of $Z$, where $Z$ and $T$ are $S$-schemes.)
\end{definition}


\begin{definition}[Schematic quotient]\label{nhq-schematic-quotient}\source{nitsure-hilbert-quot:page-0031}\uses{nhq-schematic-equivalence-relation}
We will say that a morphism $q : X\to Q$ of $S$-schemes
is a \textbf{quotient} for a
schematic equivalence relation $f: R \to X\times_SX$ over $S$ if $q$
is a \textbf{co-equaliser} for the component morphisms
$f_1,f_2 : R  \stackrel{\to}{\scriptstyle\to} X$ of
$f: R \to X\times_SX$. This means $q \circ f_1 = q\circ f_2$,
and given any $S$-scheme $Z$ and an $S$-morphism $g : X \to Z$ such that
$g\circ f_1 = g\circ f_2$, there exists a unique $S$-morphism
$h : Q \to Z$ such that $g = h\circ q$.
A schematic quotient $q: X\to Q$, when it exists, is
unique up to a unique isomorphism.
\end{definition}

\begin{claim}[Functor-of-points caution]\label{nhq-quotient-functor-points-caution}\source{nitsure-hilbert-quot:page-0031}\uses{nhq-schematic-quotient}
Even when $q:X\to Q$ is a schematic quotient for $R$, the map
$q(T):X(T)\to Q(T)$ need not be a quotient of the equivalence relation
$R(T)$ in sets: it may fail to be surjective and may identify distinct
equivalence classes.
\end{claim}

\begin{definition}[Effective schematic quotient]\label{nhq-effective-schematic-quotient}\source{nitsure-hilbert-quot:page-0031}\uses{nhq-schematic-quotient}
The quotient $q:X\to Q$ is called \textbf{effective} if the induced
morphism $(f_1,f_2):R\to X\times_QX$ is an isomorphism of $S$-schemes.
\end{definition}

But under the powerful assumption of projectivity,
Grothendieck proved an existence result for quotients, to which we
devote the rest of this section.

We will need the following elementary lemma from Grothendieck's
theory of faithfully flat descent (this is a special case of
[SGA 1] Expos\'e VIII Corollary 1.9).
The reader can consult the lectures of Vistoli
[V] for an exposition of descent.

\begin{lemma}\label{nhq-descent-subshemes}\dcref{6.7}\source{nitsure-hilbert-quot:page-0032}
(1) Any faithfully flat quasi-compact morphism of schemes $f : X \to Y$
is an effective epimorphism, that is, $f$ is a co-equaliser for
the projections $p_1,p_2: X\times_YX   \stackrel{\to}{\scriptstyle\to} X$.

(2) Let $p : D\to H$ be a faithfully flat quasi-compact morphism.
Let $Z \subset D$ be a closed subscheme
such that
$$p_1^{-1}Z = p_2^{-1}Z \subset D\times_HD$$
where $p_1,p_2 :  D\times_HD \stackrel{\to}{\scriptstyle\to} D$
are the projections, and $p_i^{-1}Z$ is the schematic
inverse image of $Z$ under $p_i$.
Then there exists a unique closed subscheme $Q$ of $H$
such that $Z = p^{-1}Q \subset D$. By base-change from $p : D\to H$,
it follows that the induced morphism $p|_Z : Z\to Q$
is faithfully flat and quasi-compact.
\end{lemma}

\begin{proof}
\notready The source states this faithfully flat descent lemma and appends a square without giving a proof.
\end{proof}


The idea of using Hilbert schemes to make
quotients of flat projective equivalence relations is due to Grothendieck,
who used it in his construction of a relative Picard scheme.
In set-theoretic terms, the idea is actually very simple:
Let $X$ be a set,
and $R\subset X\times X$ an equivalence relation on $X$.
Let $H$ be the power set of $X$ (means the set of all subsets of $X$),
and let ${\varphi}: X \to H$ be the map which sends $x\in X$ to its
equivalence class $[x] \in H$. If $Q\subset H$ is the image of ${\varphi}$,
then the induced map $q: X\to Q$ is the quotient of $X$ modulo $R$ in the
category of sets. The scheme-theoretic analogue of the above is the
following theorem of Grothendieck, where the Hilbert scheme of $X$
plays the role of power set. The first detailed proof appeared in
Altman and Kleiman [A-K 2].


\begin{theorem}\label{nhq-projective-flat-quotient}\dcref{6.8}\source{nitsure-hilbert-quot:page-0033}\uses{nhq-quasi-projective-quot,nhq-morphism-scheme,nhq-schematic-equivalence-relation,nhq-schematic-quotient,nhq-effective-schematic-quotient,nhq-descent-subshemes,nhq-hilbert-functor}
Let $S$ be a noetherian scheme, and let
$X\to S$ be a quasi-projective morphism.
Let $f : R \to X\times_SX$ be a schematic
equivalence relation on $X$ over $S$, such that
the projections $f_1,f_2 : R \stackrel{\to}{\scriptstyle\to} X$
are proper and flat.
Then a schematic quotient $X\to Q$ exists over $S$.
Moreover, $Q$ is quasi-projective over $S$,
the morphism $X \to Q$ is faithfully flat and projective, and
the induced morphism $(f_1,f_2) : R\to X\times_QX$ is an isomorphism (the
quotient is effective).
\end{theorem}


\begin{proof} (Following Altman and Kleiman [A-K 2])
The properness of $f_i$
together with separatedness of $X\to S$ implies properness of
$f : R \to X\times_SX$.
Also, $f$ is functorially injective by definition
of a schematic equivalence relation.
It follows that $f$ is a closed embedding, which
allows us to regard $R$ as a closed subscheme of $X\times_SX$.
This defines an element $(R)$ of $\Hilb_{X/S}(X)$, as the projection
$p_2|_R = f_2$ is proper and flat.




By Theorem \ref{nhq-quasi-projective-quot},
there exists a scheme $\hilb_{X/S}$ which represents the
functor $\Hilb_{X/S}$.
As the parameter scheme $X$ is noetherian and as the Hilbert polynomial
is locally constant,
only finitely many polynomials $\Phi$ occur as Hilbert
polynomials of fibers of $f_2: R \to X$ with respect to a
chosen relatively very ample line bundle $L$ on $X$ over $S$.
Let $H$ be the finite disjoint union of the corresponding
open subschemes
$\hilb_{X/S}^{\Phi,L}$ of $\hilb_{X/S}$.
Then $H$ is a quasi-projective scheme over $S$ as
each $\hilb_{X/S}^{\Phi,L}$ is so by
Theorem \ref{nhq-quasi-projective-quot}.
The family $(R) \in  \Hilb_{X/S}(X)$
therefore defines a classifying morphism
${\varphi}: X \to H$, with the property that $\varphi^*D = R$
where $D\subset X\times_SH$ denotes the restriction to $H$
of the universal family over $\hilb_{X/S}$,
and $\varphi^*D$ denotes $(\id_X\times\,{\varphi})^{-1} D$.
Also, note that the projection $p: D\to H$ is proper and flat.
If $X$ is non-empty then each fiber of $f_2 : R \to X$
is also non-empty as the diagonal $\Delta_X$ is contained in $R$,
and therefore the Hilbert polynomial of each fiber is non-zero.
Hence $p: D\to H$ is surjective, and so $p$ is faithfully flat.


For any $S$-scheme $T$, it follows from its definition that
a $T$-valued point of $D$ is a pair $(x,V)$ with
$x\in X(T)$ and $V \in H(T)$ such that
$$x\in V$$
where the notation ``$x\in V$'' more precisely means that
the graph morphism $(x, \id_T) : T \to X\times_ST$
factors via $V \subset X\times_ST$. With this notation,
we will establish the following crucial property:


\medskip



\textbf{(***)} \textit{For any $S$-scheme $T$, and
$T$-valued points $x,y \in X(T)$,
the following equivalences hold:}
$(x,y) \in R(T) \Leftrightarrow x \in \varphi(y)
\Leftrightarrow \varphi(x) = \varphi(y) \in H(T)$.



\medskip


For this, note that for any $x,y\in X(T)$, the morphism
$(x,y): T \to X\times_SX$  factors as the composite
$$T \stackrel{(x,\id_T)}{\to}X\times_S T
\stackrel{\id_X \times\, y}{\to}  X\times_S X$$
As $\varphi^*D = R$ and $(\varphi\circ y)^*D = \varphi(y)$,
it follows that $y^*R =  \varphi(y)$, in other words,
the schematic inverse image of $R\subset X\times_S X$ under
$\id_X\times\, y : X\times_S T \to X\times_S X$
is $\varphi(y)\subset X\times_S T$. Hence the above factorisation
of $(x,y) : T\to X\times_SX$ shows that
\begin{eqnarray*}
&& (x,y) \in R(T) \\
& \Leftrightarrow
& \mbox{~the~morphism~} (x,\id_T):T \to X\times_ST
  \mbox{~factors~via~} \varphi(y)\subset X\times_ST \\
& \Leftrightarrow & x \in \varphi(y)
\end{eqnarray*}
Moreover, $x \in \varphi(x)$ as $\Delta_X \subset R$.
Therefore, if $\varphi(x) = \varphi(y)$ then
$x \in \varphi(y)$.

It now only remains to prove that
if $(x,y) \in R(T)$ then $\varphi(x) = \varphi(y)$,
that is, the subschemes
$\varphi(x)$ and $\varphi(y)$ of $X\times_S T$ are identical.
Note that
$\varphi(x) = (\varphi\circ x)^*D = x^*\varphi^*D =
x^*R$ and similarly $\varphi(y) = y^*R$, therefore
we wish to show that $x^*R  = y^*R$. To show this
in terms of functor of points, for any $T$-scheme $u: U \to T$
we just have to show that
$(x^*R)(U) = (y^*R)(U)$ as subsets of $(X\times_ST)(U)$.
As $x^*R$ is the inverse image of $R$ under
$\id_X\times\, x : X\times_ST \to X\times_SX$, it follows that
a $U$-valued point of $x^*R$ is the same as
an element $z\in X(U)$
such that $(\id_X\times\, x)\circ (z,u) \in R(U)$.
But as $(\id_X\times\, x)\circ (z,u) = (z, x\circ u)$, it follows that
$$z \in (x^*R)(U) \Leftrightarrow (z, x\circ u) \in R(U)$$
As $R(U)$ is an equivalence relation on the set $X(U)$, and
as by assumption $(x,y) \in R(T)$, we have
$(x\circ u, y\circ u) \in R(U)$, and so by transitivity
we have
$$z \in (x^*R)(U) \Leftrightarrow (z, x\circ u) \in R(U)
\Leftrightarrow (z, y\circ u) \in R(U)
\Leftrightarrow z \in (y^*R)(U)$$
Hence the subschemes $x^*R$ and $y^*R$ of $X\times_ST$
have the same $U$-valued points
for any $T$-scheme $U$, and therefore $x^*R = y^*R$, as was to be shown.
This completes the proof of the assertion \textbf{(***)}.






The graph morphism $(\id_X, {\varphi}) : X \to X\times_S H$
is a closed embedding as $H$ is separated over $S$.
As $\Delta_X \subset R$ and as
${\varphi}^* D = R$, it follows that
$(\id_X, {\varphi})$ factors through $D\subset X\times_SH$.
Thus, we get a closed subscheme $\Gamma_{\varphi} \subset D$,
which is the isomorphic image of $X$ under
$(\id_X, {\varphi})$. We wish to apply the
Lemma \ref{nhq-descent-subshemes}.(2)
to the faithfully flat quasi-compact morphism
$p: D \to H$ and the closed subscheme $Z = \Gamma_{\varphi} \subset D$.

Any $T$-valued point of
$\Gamma_{\varphi}$ is a pair $(x, {\varphi}(x)) \in D(T)$ where
$x\in X(T)$. Any $T$-valued point of $D\times_HD$ is a triple
$(x,y,V)$ where $x,y\in X(T)$ and $V \in H(T)$
such that $x, y \in V$.
Under the projections
$p_1,p_2 :  D\times_HD \stackrel{\to}{\scriptstyle\to} D$,
we have $p_1(x,y,V) = (x,V)$ and $p_2(x,y,V) = (y,V)$.
We now have
\begin{eqnarray*}
&& p_1(x,y,V) \in \Gamma_{\varphi}(T)\\
& \Leftrightarrow & (x,V) \in \Gamma_{\varphi}(T) \mbox{ and }y\in V \\
& \Leftrightarrow & V = \varphi(x)\mbox{ and }y\in V \\
& \Leftrightarrow & y \in \varphi(x) = V \\
& \Leftrightarrow & \varphi(y) = \varphi(x) =V ~~(\mbox{by the property
\textbf{(***)}).}
\end{eqnarray*}
Similarly, we have $p_2(x,y,V) \in \Gamma_{\varphi}(T)$ if and only if
$\varphi(y) = \varphi(x) =V$.
Therefore, $p_1(x,y,V) \in \Gamma_{\varphi}(T)$ if and only if
$p_2(x,y,V) \in \Gamma_{\varphi}(T)$.
This holds for all $T$-valued points for all $S$-schemes $T$, and so
$p_1^{-1}\Gamma_{\varphi} = p_2^{-1}\Gamma_{\varphi} \subset D\times_HD$.
Therefore by Lemma \ref{nhq-descent-subshemes}(2)
there exists a unique closed subscheme $Q \subset H$ such that
$\Gamma_{\varphi}$ is the pull-back of $Q$ under $D\to H$.
Let $p : \Gamma_{\varphi} \to Q$ be the morphism induced by the restriction
to $\Gamma_{\varphi}$ of $p: D\to H$.
Let $q: X\to Q$ be defined as the composite
$$X \stackrel{(\id_X,\varphi)}{\to} \Gamma_{\varphi}
\stackrel{p}{\to} Q$$
Then note that the composite
$X \stackrel{q}{\to} Q \hra H$ equals $\varphi$.

We will now show that $q: X\to Q$ as defined above is the desired
quotient of $X$ by $R$, with the required properties.

\textbf{(i) Quasi-projectivity of $Q\to S$ }: This is satisfied as
$Q$ is closed in $H$ and $H$ is quasi-projective over $S$.

\textbf{(ii) Faithful flatness and projectivity of $q$ }:
This follows
by base change from the faithfully flat projective morphism $p: D\to H$,
as the following square is Cartesian.
$$\begin{array}{ccc}
X & \stackrel{(\id_X,\varphi)}{\to} & D \\
{\scriptstyle q} \dna ~ & \square & ~ \dna {\scriptstyle p} \\
Q & \hra & H
\end{array}$$

\textbf{(iii) Exactness of $R \stackrel{\to}{\scriptstyle\to} X \to Q$
and the isomorphism $R \to X\times_QX$ }:
By \textbf{(***)}, for any $T$-valued points $x,y \in X(T)$, we have
$(x,y)\in R(T)$ if and only if $\varphi(x) = \varphi(y)$.
This shows that the composite $R \stackrel{f_1}{\to} X \stackrel{q}{\to} Q$
equals the composite $R \stackrel{f_2}{\to} X \stackrel{q}{\to} Q$, and
the induced morphism $(f_1,f_2) : R \to X\times_QX$
is an isomorphism, by showing these statements
hold at the level of functor of points. Under the isomorphism
$R \to X\times_QX$, the
morphisms $f_1,f_2 : R  \stackrel{\to}{\scriptstyle\to} X$
become the projection morphisms $p_1,p_2 : X\times_QX
\stackrel{\to}{\scriptstyle\to} X$.
By Lemma \ref{nhq-descent-subshemes}, the
morphism $q: X \to Q$ is a co-equaliser for $p_1,p_2$, and so $q$
is a co-equaliser for $f_1,f_2$.

This completes the proof of Theorem \ref{nhq-projective-flat-quotient}.\end{proof}



What Altman and Kleiman actually prove in [A-K 2] is
a strongly projective form of the above theorem
(without a noetherian assumption),
using the hypothesis of strong quasi-projectivity
in the following places in the above proof: if $X\to S$ is
strongly quasi-projective then $H \to S$ will again be so
by Theorem \ref{nhq-strongly-quasi-projective-quot}, and therefore
$Q$ will be strongly quasi-projective over $S$. Moreover, $D\to H$
will be strongly projective, and therefore by base-change
$X\to Q$ will be strongly projective.
In the noetherian case, this gives us the following result.



\begin{theorem}\label{nhq-strongly-projective-quotient}\dcref{6.9}\source{nitsure-hilbert-quot:page-0035}\uses{nhq-projective-flat-quotient,nhq-effective-schematic-quotient,nhq-strongly-quasi-projective-quot}
Let $S$ be a noetherian scheme, and let
$X\to S$ be a strongly quasi-projective morphism.
Let $f : R \to X\times_SX$ be a schematic
equivalence relation on $X$ over $S$, such that
the projections $f_1,f_2 : R \stackrel{\to}{\scriptstyle\to} X$
are proper and flat.
Then a schematic quotient $X\to Q$ exists over $S$.
Moreover, the quotient is effective,
the morphism $X \to Q$ is
faithfully flat and strongly projective, and $Q$ is strongly
quasi-projective over $S$.
\end{theorem}

\begin{proof}
\notready The source states this strongly projective variant as a consequence of the preceding argument and does not supply a separate proof.
\end{proof}














\bigskip
\bigskip
